% !TEX program = xelatex
\documentclass[11pt,a4paper]{article}
\usepackage[margin=2.2cm,top=2cm]{geometry}
\usepackage{fontspec}
\usepackage{polyglossia}
\setdefaultlanguage{turkish}
\usepackage{amsmath,amssymb}
\usepackage{graphicx}
\usepackage{booktabs}
\usepackage{enumitem}
\usepackage{multicol}
\usepackage[most]{tcolorbox}
\usepackage{fancyhdr}
\usepackage{titlesec}
\usepackage{xcolor}
\usepackage{needspace}
\usepackage[hidelinks]{hyperref}

\definecolor{anamavi}{RGB}{0,45,98}
\definecolor{cozumyesil}{RGB}{0,110,50}
\definecolor{uyarikirmizi}{RGB}{180,20,20}

\titleformat{\section}{\large\bfseries\color{anamavi}}{\thesection.}{0.5em}{}
\titleformat{\subsection}{\normalsize\bfseries\color{anamavi}}{\thesubsection}{0.5em}{}

\pagestyle{fancy}\fancyhf{}
\lhead{\small MMAT105 Genel Matematik}
\rhead{\small Hafta 10 -- Ders Notu}
\cfoot{\small\thepage}
\renewcommand{\headrulewidth}{0.4pt}

\newtcolorbox{kural}[1][Kural]{colback=anamavi!5,colframe=anamavi,title=#1,fonttitle=\bfseries,breakable,left=4pt,right=4pt,top=2pt,bottom=2pt}
\newtcolorbox{dikkat}[1][Dikkat]{colback=uyarikirmizi!5,colframe=uyarikirmizi,title=#1,fonttitle=\bfseries,breakable,left=4pt,right=4pt,top=2pt,bottom=2pt}
\newcounter{ornek}
\newenvironment{ornek}{\refstepcounter{ornek}\par\medskip\noindent\textbf{\color{anamavi}Örnek \theornek.}\ }{\par}
\newenvironment{cozum}{\begin{tcolorbox}[colback=cozumyesil!6,colframe=cozumyesil,boxrule=0.5pt,left=4pt,right=4pt,top=2pt,bottom=2pt,breakable]\textbf{\color{cozumyesil}Çözüm.}\ }{\end{tcolorbox}}
\newcounter{alis}
\newenvironment{alistirma}{\refstepcounter{alis}\par\medskip\needspace{9\baselineskip}\noindent\textbf{\color{anamavi}Alıştırma \thealis.}\ }{\par}
\newcommand{\kendinizoncesi}{\needspace{12\baselineskip}}
\newcommand{\ab}[1]{\overline{#1}}
\setlength{\parindent}{0pt}
\setlength{\parskip}{3pt}

\begin{document}

% ---------------- Başlık ----------------
\begin{minipage}[c]{\textwidth}
{\LARGE\bfseries\color{anamavi} MMAT105 Genel Matematik}\\[4pt]
{\large Hafta 10 Ders Notu: Köklü Sayılar\\ (Özellikler, Dört İşlem, Eşlenik)}\\[6pt]
{\small 2026--2027 Güz}
\end{minipage}

\vspace{4pt}\noindent\textcolor{anamavi}{\rule{\textwidth}{1pt}}

\textbf{Bu notun kullanımı.} Her bölümde önce kısa bir kural kutusu, ardından çözümlü örnekler yer alır. Örnekleri çözümüne bakmadan önce kendiniz deneyin. Notun sonunda önce çözümlü alıştırmalar, ardından cevap anahtarlı ``Kendiniz Çözün'' alıştırmaları yer alır.

% =====================================================
\section{Köklü Sayının Tanımı}

\begin{kural}[Tanım (\(n\ge 2\) tam sayı: kök derecesi)]
\begin{itemize}[leftmargin=*,itemsep=1pt]
\item \(a\ge 0\) için \(\sqrt{a}\), karesi \(a\) olan \textbf{negatif olmayan} sayıdır: \(\sqrt{49}=7\), \(\sqrt{0}=0\).
\item \(n\) \textbf{çift} ise \(\sqrt[n]{a}\) yalnız \(a\ge 0\) için reel sayıdır ve \(\sqrt[n]{a}\ge 0\)'dır: \(\sqrt[4]{81}=3\).
\item \(n\) \textbf{tek} ise \(\sqrt[n]{a}\) her reel \(a\) için tanımlıdır ve işareti \(a\)'nın işaretidir: \(\sqrt[3]{-8}=-2\), \(\sqrt[5]{-1}=-1\).
\end{itemize}
\textbf{Tam kareler:} \(1,4,9,16,25,36,49,64,81,100,121,144,169,196,225,256,289,324,361,400\)\\
\textbf{Tam küpler:} \(1,8,27,64,125,216,343,512,729,1000\)
\end{kural}

\begin{dikkat}
\(\sqrt{9}=3\)'tür, \(\pm 3\) değildir. \(\sqrt{-4}\), \(\sqrt[4]{-16}\) reel sayı değildir. \(\sqrt{a+b}\neq\sqrt a+\sqrt b\): \(\sqrt{9+16}=5\) ama \(\sqrt9+\sqrt{16}=7\).
\end{dikkat}

\begin{ornek}
\(\sqrt{0{,}09}+\sqrt[3]{-0{,}008}\) ve \(\sqrt[4]{\dfrac{1}{81}}\) değerleri kaçtır?
\end{ornek}
\begin{cozum}
\(\sqrt{0{,}09}=0{,}3\) \ (\(0{,}3^2=0{,}09\)); \ \(\sqrt[3]{-0{,}008}=-0{,}2\) \ (\((-0{,}2)^3=-0{,}008\)). Toplam: \(0{,}3-0{,}2=0{,}1\).\\
\(\sqrt[4]{\frac{1}{81}}=\frac13\) \ (\(3^4=81\)).
\end{cozum}

\begin{ornek}
\(\sqrt[4]{10-2x}\) bir reel sayı belirtiyorsa \(x\)'in alabileceği en büyük tam sayı değeri kaçtır? \(x\) kaç farklı pozitif tam sayı değeri alabilir?
\end{ornek}
\begin{cozum}
Çift dereceli kök: \(10-2x\ge 0\Rightarrow x\le 5\). En büyük tam sayı değeri \(5\). Pozitif tam sayı değerleri \(1,2,3,4,5\): \(5\) tane.
\end{cozum}

% =====================================================
\section{Kesirli Üs ve \texorpdfstring{\(\sqrt{a^2}=|a|\)}{karekok a kare}}

\begin{kural}[Kural]
\(a>0\), \(m\) tam sayı ve \(n\ge 2\) tam sayı ise \quad \(a^{\frac mn}=\sqrt[n]{a^m}=\left(\sqrt[n]{a}\right)^m\). \quad Örnek: \(\sqrt[3]{x^2}=x^{\frac23}\), \ \(\dfrac{1}{\sqrt x}=x^{-\frac12}\) (\(x>0\)).\\[2pt]
Her reel \(a\) için \(\sqrt{a^2}=|a|\); \quad \(\sqrt[n]{a^n}=a\) (\(n\) tek), \quad \(\sqrt[n]{a^n}=|a|\) (\(n\) çift).
\end{kural}

\begin{dikkat}
Kesirli üste pay üs, payda kök derecesidir: \(a^{\frac23}=\sqrt[3]{a^2}\). \(\sqrt{x^2}=x\) eşitliği yalnız \(x\ge 0\) iken doğrudur; genel olarak \(\sqrt{x^2}=|x|\) (mutlak değer: 7.~hafta).
\end{dikkat}

\begin{ornek}
\(32^{\frac35}+81^{-\frac14}\) işleminin sonucu kaçtır?
\end{ornek}
\begin{cozum}
\(32^{\frac35}=\left(\sqrt[5]{32}\right)^3=2^3=8\); \quad \(81^{-\frac14}=\dfrac{1}{81^{\frac14}}=\dfrac{1}{\sqrt[4]{81}}=\dfrac13\). \quad Sonuç: \(8+\dfrac13=\dfrac{25}{3}\).
\end{cozum}

\begin{ornek}
\(\sqrt[3]{4^x}=16\) olduğuna göre \(x\) kaçtır?
\end{ornek}
\begin{cozum}
\(\sqrt[3]{4^x}=\sqrt[3]{2^{2x}}=2^{\frac{2x}{3}}\) ve \(16=2^4\). \quad \(\dfrac{2x}{3}=4\Rightarrow x=6\). \quad Kontrol: \(\sqrt[3]{4^6}=\sqrt[3]{4096}=16\).
\end{cozum}

\begin{ornek}
\(\sqrt{(\sqrt3-2)^2}+\sqrt{(\sqrt3-1)^2}\) işleminin sonucu kaçtır?
\end{ornek}
\begin{cozum}
\(1<\sqrt3<2\) (çünkü \(1<3<4\)). \ \(\sqrt3-2<0\Rightarrow|\sqrt3-2|=2-\sqrt3\); \ \(\sqrt3-1>0\Rightarrow|\sqrt3-1|=\sqrt3-1\).
\[(2-\sqrt3)+(\sqrt3-1)=1\]
\end{cozum}

\begin{ornek}
\(x<0<y\) olmak üzere \(\sqrt{x^2}+\sqrt[3]{y^3}-\sqrt{(x-y)^2}\) ifadesinin eşiti nedir?
\end{ornek}
\begin{cozum}
\(\sqrt{x^2}=|x|=-x\) (\(x<0\)); \quad \(\sqrt[3]{y^3}=y\) (tek derece); \quad \(x-y<0\Rightarrow\sqrt{(x-y)^2}=|x-y|=y-x\).
\[-x+y-(y-x)=-x+y-y+x=0\]
\end{cozum}

% =====================================================
\section{Kök Dışına Çıkarma, Kök İçine Alma, İki Tam Sayı Arası}

\begin{kural}[Kural]
\(a,b\ge 0\) için \(\sqrt{a^2 b}=a\sqrt b\) \ ve \ \(\sqrt[n]{a^n b}=a\sqrt[n]{b}\). Kök dışına çıkarırken içerideki en büyük tam kare (tam küp) çarpan ayrılır; kök içine alırken dıştaki sayının \(n\). kuvveti içeri girer.\\
\(k^2<a<(k+1)^2\) ise \(k<\sqrt a<k+1\).
\end{kural}

\begin{ornek}
\(\sqrt{2^5\cdot 3^3}\) ve \(\sqrt[3]{2^7\cdot 5^4}\) sayılarını kök dışına çıkararak sadeleştiriniz.
\end{ornek}
\begin{cozum}
Üsleri kök derecesinin katlarına ayıralım:
\(\sqrt{2^4\cdot 2\cdot 3^2\cdot 3}=2^2\cdot 3\sqrt{2\cdot 3}=12\sqrt6\); \quad
\(\sqrt[3]{2^6\cdot 2\cdot 5^3\cdot 5}=2^2\cdot 5\sqrt[3]{2\cdot 5}=20\sqrt[3]{10}\).
\end{cozum}

\begin{ornek}
\(a\sqrt5=\sqrt{180}\) ve \(2\sqrt[3]{5}=\sqrt[3]{b}\) olduğuna göre \(a+b\) kaçtır?
\end{ornek}
\begin{cozum}
\(\sqrt{180}=\sqrt{36\cdot 5}=6\sqrt5\Rightarrow a=6\). \quad \(2\sqrt[3]{5}=\sqrt[3]{2^3\cdot 5}=\sqrt[3]{40}\Rightarrow b=40\). \quad \(a+b=46\).
\end{cozum}

\begin{ornek}
\(\sqrt{130}\), \(\sqrt[3]{100}\) ve \(2\sqrt7\) sayıları hangi iki ardışık tam sayı arasındadır?
\end{ornek}
\begin{cozum}
\(121<130<144\Rightarrow 11<\sqrt{130}<12\). \quad \(64<100<125\Rightarrow 4<\sqrt[3]{100}<5\).\\
\(2\sqrt7=\sqrt{28}\), \ \(25<28<36\Rightarrow 5<2\sqrt7<6\).
\end{cozum}

% =====================================================
\section{Toplama, Çıkarma, Çarpma ve Bölme}

\begin{kural}[Dört işlem]
\begin{itemize}[leftmargin=*,itemsep=1pt]
\item Yalnız \textbf{benzer kökler} (aynı derece, aynı iç) toplanır: \(a\sqrt[n]{x}+b\sqrt[n]{x}=(a+b)\sqrt[n]{x}\). Önce sadeleştirin.
\item Aynı dereceli kökler (\(n\) çift ise \(a,b\ge 0\)): \(\sqrt[n]{a}\cdot\sqrt[n]{b}=\sqrt[n]{ab}\), \ \(\dfrac{\sqrt[n]{a}}{\sqrt[n]{b}}=\sqrt[n]{\dfrac ab}\) (\(b\neq 0\)).
\item Derece genişletme (\(a\ge 0\)): \(\sqrt[n]{a^m}=\sqrt[nk]{a^{mk}}\). Farklı dereceler EKOK'ta eşitlenir.
\item Kökün kökü (\(a\ge 0\)): \(\sqrt[m]{\sqrt[n]{a}}=\sqrt[mn]{a}\).
\end{itemize}
\end{kural}

\begin{dikkat}[Sık yapılan hatalar]
\(\sqrt2+\sqrt3\neq\sqrt5\). \ \(\sqrt{-4}\cdot\sqrt{-9}=\sqrt{36}\) yazılamaz (\(\sqrt{-4}\) reel değil). \ \(\sqrt2\cdot\sqrt[3]{3}\neq\sqrt[6]{6}\): dereceyi genişletirken içerinin üssü de çarpılır.
\end{dikkat}

\begin{ornek}
\(3\sqrt8-\sqrt{50}+\sqrt{32}\) işleminin sonucu nedir?
\end{ornek}
\begin{cozum}
\(3\sqrt8=3\cdot 2\sqrt2=6\sqrt2\), \ \(\sqrt{50}=5\sqrt2\), \ \(\sqrt{32}=4\sqrt2\). \quad \(6\sqrt2-5\sqrt2+4\sqrt2=5\sqrt2\).
\end{cozum}

\begin{ornek}
\(2\sqrt[3]{24}-\sqrt[3]{81}\) işleminin sonucu nedir?
\end{ornek}
\begin{cozum}
\(\sqrt[3]{24}=\sqrt[3]{8\cdot 3}=2\sqrt[3]{3}\), \ \(\sqrt[3]{81}=\sqrt[3]{27\cdot 3}=3\sqrt[3]{3}\). \quad \(2\cdot 2\sqrt[3]{3}-3\sqrt[3]{3}=\sqrt[3]{3}\).
\end{cozum}

\begin{ornek}
\((\sqrt{12}+\sqrt{75})\cdot\sqrt3\) işleminin sonucu kaçtır?
\end{ornek}
\begin{cozum}
Dağılma özelliği: \(\sqrt{12}\cdot\sqrt3+\sqrt{75}\cdot\sqrt3=\sqrt{36}+\sqrt{225}=6+15=21\).\\
(Diğer yol: \(\sqrt{12}+\sqrt{75}=2\sqrt3+5\sqrt3=7\sqrt3\), \ \(7\sqrt3\cdot\sqrt3=7\cdot 3=21\).)
\end{cozum}

\begin{ornek}
\(\sqrt2\cdot\sqrt{18}+\dfrac{\sqrt{63}}{\sqrt7}\) işleminin sonucu kaçtır?
\end{ornek}
\begin{cozum}
\(\sqrt2\cdot\sqrt{18}=\sqrt{36}=6\), \quad \(\dfrac{\sqrt{63}}{\sqrt7}=\sqrt9=3\). \quad Sonuç: \(9\).
\end{cozum}

\begin{ornek}
\(\sqrt[3]{2}\cdot\sqrt[4]{2}\) çarpımını tek kök olarak yazınız. \(\sqrt{3\sqrt[3]{3}}\) ifadesinin en sade biçimi nedir?
\end{ornek}
\begin{cozum}
EKOK\((3,4)=12\): \(\sqrt[3]{2}=\sqrt[12]{2^4}\), \(\sqrt[4]{2}=\sqrt[12]{2^3}\) \(\Rightarrow\) \(\sqrt[12]{2^7}=\sqrt[12]{128}\).\\
\(3\)'ü içteki köke alalım: \(3\sqrt[3]{3}=\sqrt[3]{27\cdot 3}=\sqrt[3]{81}\). \quad \(\sqrt{\sqrt[3]{81}}=\sqrt[6]{81}=\sqrt[6]{3^4}=\sqrt[3]{3^2}=\sqrt[3]{9}\).
\end{cozum}

\begin{ornek}
\(a=\sqrt2\), \(b=\sqrt[3]{3}\), \(c=\sqrt[4]{5}\) sayılarını küçükten büyüğe sıralayınız.
\end{ornek}
\begin{cozum}
EKOK\((2,3,4)=12\): \ \(a=\sqrt[12]{2^6}=\sqrt[12]{64}\), \ \(b=\sqrt[12]{3^4}=\sqrt[12]{81}\), \ \(c=\sqrt[12]{5^3}=\sqrt[12]{125}\). \quad \(64<81<125\Rightarrow a<b<c\).
\end{cozum}

\begin{ornek}
\(\sqrt5\approx 2{,}236\) olduğuna göre \(\sqrt{20}+\sqrt{45}-\sqrt{80}\) ifadesinin yaklaşık değeri kaçtır?
\end{ornek}
\begin{cozum}
\(\sqrt{20}=2\sqrt5\), \ \(\sqrt{45}=3\sqrt5\), \ \(\sqrt{80}=4\sqrt5\). \quad \(2\sqrt5+3\sqrt5-4\sqrt5=\sqrt5\approx 2{,}236\).
\end{cozum}

% =====================================================
\section{Paydayı Rasyonel Yapma ve Eşlenik}

\begin{kural}[Kural]
\begin{itemize}[leftmargin=*,itemsep=1pt]
\item Tek terimli payda (\(b>0\)): \(\dfrac{a}{\sqrt b}=\dfrac{a\sqrt b}{b}\). \ Küp kökte: \(\sqrt[3]{x}\cdot\sqrt[3]{x^2}=x\).
\item \textbf{İki kare farkı} (yalnız eşlenik çarpımı için kısa kural): \((x-y)(x+y)=x^2-y^2\), \ yani \((\sqrt a-\sqrt b)(\sqrt a+\sqrt b)=a-b\) (\(a,b\ge 0\)).
\item Payda \(\sqrt a\pm\sqrt b\) veya \(c\pm\sqrt b\) ise pay ve payda, paydanın \textbf{eşleniği} (\(\sqrt a\mp\sqrt b\), \(c\mp\sqrt b\)) ile çarpılır.
\end{itemize}
\end{kural}

\begin{dikkat}
Pay ve payda \textbf{aynı} ifadeyle çarpılmalı. \((\sqrt3+1)^2\) gibi bir kareyi dağılma özelliğiyle adım adım açın: \((\sqrt3+1)(\sqrt3+1)=3+\sqrt3+\sqrt3+1=4+2\sqrt3\) (\(3+1=4\) değil!).
\end{dikkat}

\begin{ornek}
\(\dfrac{15}{\sqrt5}+\sqrt{20}\) işleminin sonucu nedir?
\end{ornek}
\begin{cozum}
\(\dfrac{15}{\sqrt5}=\dfrac{15\sqrt5}{5}=3\sqrt5\), \ \(\sqrt{20}=2\sqrt5\). \quad \(3\sqrt5+2\sqrt5=5\sqrt5\).
\end{cozum}

\begin{ornek}
\(\dfrac{2}{\sqrt[3]{4}}\) kesrinin paydasını rasyonel yapınız.
\end{ornek}
\begin{cozum}
\(\sqrt[3]{4}\cdot\sqrt[3]{2}=\sqrt[3]{8}=2\) olduğundan pay ve payda \(\sqrt[3]{2}\) ile çarpılır:
\(\dfrac{2\sqrt[3]{2}}{\sqrt[3]{8}}=\dfrac{2\sqrt[3]{2}}{2}=\sqrt[3]{2}\).
\end{cozum}

\begin{ornek}
\(\dfrac{3}{\sqrt7+2}\) kesrinin paydasını rasyonel yapınız.
\end{ornek}
\begin{cozum}
Eşlenik \(\sqrt7-2\): \ payda \((\sqrt7+2)(\sqrt7-2)=7-4=3\). \quad \(\dfrac{3(\sqrt7-2)}{3}=\sqrt7-2\).
\end{cozum}

\begin{ornek}
\(\dfrac{\sqrt5+\sqrt3}{\sqrt5-\sqrt3}\) kesrini sadeleştiriniz.
\end{ornek}
\begin{cozum}
Pay ve payda \(\sqrt5+\sqrt3\) ile çarpılır. Payda: \(5-3=2\).\\
Pay (dağılma): \((\sqrt5+\sqrt3)(\sqrt5+\sqrt3)=5+\sqrt{15}+\sqrt{15}+3=8+2\sqrt{15}\). \quad \(\dfrac{8+2\sqrt{15}}{2}=4+\sqrt{15}\).
\end{cozum}

\begin{ornek}
\(x=\dfrac{1}{\sqrt2-1}\) olduğuna göre \(x-\dfrac1x\) kaçtır?
\end{ornek}
\begin{cozum}
\(x=\dfrac{\sqrt2+1}{(\sqrt2-1)(\sqrt2+1)}=\dfrac{\sqrt2+1}{2-1}=\sqrt2+1\); \quad \(\dfrac1x=\sqrt2-1\).
\quad \(x-\dfrac1x=(\sqrt2+1)-(\sqrt2-1)=2\).
\end{cozum}

% =====================================================
\section{İç İçe Kökler}

\begin{kural}[Kural]
\(x\ge y\ge 0\), \ \(x+y=a\) \ ve \ \(x\cdot y=b\) ise
\[\sqrt{a+2\sqrt b}=\sqrt x+\sqrt y,\qquad \sqrt{a-2\sqrt b}=\sqrt x-\sqrt y.\]
İçteki kökün önünde \(2\) yoksa önce oluşturun: \(\sqrt{60}=2\sqrt{15}\), \ \(\sqrt{56}=2\sqrt{14}\).
\end{kural}

\begin{ornek}
\(\sqrt{11+2\sqrt{30}}\) ve \(\sqrt{9-\sqrt{56}}\) ifadelerinin eşitleri nedir?
\end{ornek}
\begin{cozum}
Toplamı \(11\), çarpımı \(30\): \(6\) ve \(5\) \(\Rightarrow\) \(\sqrt{11+2\sqrt{30}}=\sqrt6+\sqrt5\).\\
\(\sqrt{56}=2\sqrt{14}\); toplamı \(9\), çarpımı \(14\): \(7\) ve \(2\) \(\Rightarrow\) \(\sqrt{9-2\sqrt{14}}=\sqrt7-\sqrt2\) (büyük olan önce).
\end{cozum}

\begin{ornek}\textit{(İleri örnek)}
\(x=\sqrt{20+\sqrt{20+\sqrt{20+\cdots}}}\) olduğuna göre \(x\) kaçtır?
\end{ornek}
\begin{cozum}
Kökün içinde yine \(x\) vardır: \(x=\sqrt{20+x}\), \(x>0\). Kare alınırsa \(x^2=20+x\Rightarrow x(x-1)=20\). Çarpımı \(20\) olan ardışık iki pozitif sayı \(5\) ve \(4\) \(\Rightarrow\) \(x=5\). (\(x>1\) iken \(x(x-1)\) büyüdükçe büyür, \(0<x\le 1\) iken \(x(x-1)\le 0\); başka pozitif çözüm yoktur.) Kontrol: \(\sqrt{20+5}=5\).
\end{cozum}

% =====================================================
\section{Özet: Formül Kartı}
\begin{center}
\renewcommand{\arraystretch}{1.6}
\begin{tabular}{@{}ll@{}}
\toprule
Konu & Formül / Kural\\\midrule
Tanım & \(n\) çift: \(a\ge 0\) olmalı, \(\sqrt[n]{a}\ge 0\); \ \(n\) tek: her \(a\), \(\sqrt[3]{-8}=-2\)\\
Kesirli üs & \(a^{\frac mn}=\sqrt[n]{a^m}=\left(\sqrt[n]{a}\right)^m\) \ (\(a>0\))\\
Mutlak değer & \(\sqrt{a^2}=|a|\); \ \(\sqrt[n]{a^n}=a\) (\(n\) tek), \ \(|a|\) (\(n\) çift)\\
Kök dışı / içi & \(\sqrt{a^2b}=a\sqrt b\), \ \(\sqrt[n]{a^nb}=a\sqrt[n]{b}\) \ (\(a,b\ge 0\))\\
Toplama & \(a\sqrt[n]{x}\pm b\sqrt[n]{x}=(a\pm b)\sqrt[n]{x}\) \ (yalnız benzer kökler)\\
Çarpma / bölme & \(\sqrt[n]{a}\sqrt[n]{b}=\sqrt[n]{ab}\), \ \(\sqrt[n]{a}/\sqrt[n]{b}=\sqrt[n]{a/b}\) \ (\(n\) çift ise \(a,b\ge 0\))\\
Ortak derece & \(\sqrt[n]{a^m}=\sqrt[nk]{a^{mk}}\); \ \(\sqrt[m]{\sqrt[n]{a}}=\sqrt[mn]{a}\) \ (\(a\ge 0\))\\
Rasyonel payda & \(\dfrac{a}{\sqrt b}=\dfrac{a\sqrt b}{b}\); \ \((\sqrt a-\sqrt b)(\sqrt a+\sqrt b)=a-b\)\\
İç içe kök & \(\sqrt{a\pm 2\sqrt b}=\sqrt x\pm\sqrt y\) \ (\(x+y=a\), \(xy=b\), \(x\ge y\))\\\bottomrule
\end{tabular}
\end{center}

% =====================================================
\newpage
% Hafta 10 -- alıştırma gövdesi (ders notu ve ayrı alıştırma kağıdı tarafından \input edilir)
\section{Çözümlü Alıştırmalar}
Her soruyu önce kendiniz deneyin, sonra çözümü okuyun.

\begin{alistirma}
\(\sqrt[3]{-8}+\sqrt[4]{16}-\sqrt[5]{-32}\) işleminin sonucu kaçtır?
\end{alistirma}
\begin{cozum}
Tek dereceli köklerde negatif içerik serbesttir:
\(\sqrt[3]{-8}=-2\) \ (\((-2)^3=-8\)), \quad \(\sqrt[4]{16}=2\) \ (\(2^4=16\)), \quad \(\sqrt[5]{-32}=-2\) \ (\((-2)^5=-32\)).
\[-2+2-(-2)=-2+2+2=2\]
\end{cozum}

\begin{alistirma}
\(\sqrt{50}+\sqrt8-\sqrt{32}\) işleminin sonucu nedir?
\end{alistirma}
\begin{cozum}
Önce her kökü sadeleştirelim (en büyük tam kare çarpanı ayıralım):\\
\(\sqrt{50}=\sqrt{25\cdot 2}=5\sqrt2\), \quad \(\sqrt8=\sqrt{4\cdot 2}=2\sqrt2\), \quad \(\sqrt{32}=\sqrt{16\cdot 2}=4\sqrt2\).
\[5\sqrt2+2\sqrt2-4\sqrt2=(5+2-4)\sqrt2=3\sqrt2\]
\end{cozum}

\begin{alistirma}
\(\dfrac{\sqrt{48}}{\sqrt3}+\dfrac{\sqrt{45}}{\sqrt5}\) işleminin sonucu kaçtır?
\end{alistirma}
\begin{cozum}
Aynı dereceli kökler bölünürken içler bölünür:
\[\frac{\sqrt{48}}{\sqrt3}=\sqrt{\frac{48}{3}}=\sqrt{16}=4,\qquad \frac{\sqrt{45}}{\sqrt5}=\sqrt{\frac{45}{5}}=\sqrt9=3,\qquad 4+3=7\]
\end{cozum}

\begin{alistirma}
\(\dfrac{4}{\sqrt5-1}\) kesrinin paydasını rasyonel yapınız.
\end{alistirma}
\begin{cozum}
Paydanın eşleniği \(\sqrt5+1\). Pay ve payda bununla çarpılır; payda iki kare farkıdır:
\((\sqrt5-1)(\sqrt5+1)=5-1=4\).
\[\frac{4(\sqrt5+1)}{4}=\sqrt5+1\]
\end{cozum}

\begin{alistirma}
\(\sqrt{2\sqrt[3]{4}}\) ifadesini tek kök olarak yazınız.
\end{alistirma}
\begin{cozum}
Önce \(2\)'yi küp köke alalım: \(2\sqrt[3]{4}=\sqrt[3]{2^3\cdot 4}=\sqrt[3]{32}\).\\
Kökün kökü: \(\sqrt{\sqrt[3]{32}}=\sqrt[2\cdot 3]{32}=\sqrt[6]{32}\). \quad (Üslü yazımla: \(2\cdot 2^{\frac23}=2^{\frac53}\), \ \(\left(2^{\frac53}\right)^{\frac12}=2^{\frac56}=\sqrt[6]{32}\).)
\end{cozum}

\begin{alistirma}
\(a=\sqrt[3]{3}\), \ \(b=\sqrt2\), \ \(c=\sqrt[6]{10}\) sayılarını küçükten büyüğe sıralayınız.
\end{alistirma}
\begin{cozum}
Kök dereceleri \(3\), \(2\), \(6\); EKOK \(=6\). Dereceleri eşitleyelim (içerinin üssü de aynı sayıyla çarpılır):
\[a=\sqrt[6]{3^2}=\sqrt[6]{9},\qquad b=\sqrt[6]{2^3}=\sqrt[6]{8},\qquad c=\sqrt[6]{10}\]
\(8<9<10\) \(\Rightarrow\) \(b<a<c\).
\end{cozum}

\begin{alistirma}
\(\sqrt{8+2\sqrt{15}}-\sqrt{8-2\sqrt{15}}\) işleminin sonucu nedir?
\end{alistirma}
\begin{cozum}
Toplamı \(8\), çarpımı \(15\) olan sayılar \(5\) ve \(3\):
\[\sqrt{8+2\sqrt{15}}=\sqrt5+\sqrt3,\qquad \sqrt{8-2\sqrt{15}}=\sqrt5-\sqrt3\quad(\text{büyük olan önce})\]
\[(\sqrt5+\sqrt3)-(\sqrt5-\sqrt3)=2\sqrt3\]
\end{cozum}

\begin{alistirma}
\(\sqrt{3x+4}=5\) \ ve \ \(\sqrt[3]{2^{y+1}}=4\) olduğuna göre \(x+y\) kaçtır?
\end{alistirma}
\begin{cozum}
1. denklem: iki tarafın karesi alınır: \(3x+4=25\Rightarrow 3x=21\Rightarrow x=7\). \ (Kontrol: \(\sqrt{25}=5\).)\\
2. denklem: kesirli üs ile \(2^{\frac{y+1}{3}}=2^2\Rightarrow\dfrac{y+1}{3}=2\Rightarrow y=5\). \ (Kontrol: \(\sqrt[3]{2^6}=\sqrt[3]{64}=4\).)
\[x+y=7+5=12\]
\end{cozum}

\kendinizoncesi
\section{Alıştırmalar (Kendiniz Çözün)}

\begin{enumerate}[label=\textbf{\arabic*.},leftmargin=*,itemsep=6pt]
\item \(\sqrt{0{,}04}+\sqrt{1{,}44}-\sqrt{0{,}25}\) işleminin sonucu kaçtır?
\item \(\sqrt[4]{81}+\sqrt[3]{-64}\) işleminin sonucu kaçtır?
\item \(\sqrt{18}\), \(\sqrt{75}\) ve \(\sqrt{128}\) sayılarını \(a\sqrt b\) biçiminde yazınız.
\item \(3\sqrt2+5\sqrt2-7\sqrt2\) işleminin sonucu nedir?
\item \(\sqrt3\cdot\sqrt6\cdot\sqrt2\) çarpımı kaçtır?
\item \(9^{\frac32}-32^{\frac25}\) işleminin sonucu kaçtır?
\item \(\sqrt{(2-\sqrt7)^2}\) ifadesinin eşiti nedir?
\item \(\sqrt{60}\) hangi iki ardışık tam sayı arasındadır?
\item \(\dfrac{6}{\sqrt3}\) kesrinin paydasını rasyonel yapınız.
\item \(\sqrt[3]{3}\cdot\sqrt3\) çarpımını tek kök olarak yazınız.
\item \(a=2\sqrt5\), \(b=3\sqrt2\), \(c=\sqrt{19}\) sayılarını küçükten büyüğe sıralayınız.
\item \(\dfrac{1}{\sqrt6-\sqrt5}\) kesrinin paydasını rasyonel yapınız.
\item \(\dfrac{2}{\sqrt3+1}+\dfrac{2}{\sqrt3-1}\) işleminin sonucu nedir?
\item \(\sqrt{6-2\sqrt5}\) ifadesinin eşiti nedir?
\item \(x=\sqrt2+1\) olduğuna göre \(x^2-2x\) kaçtır? (Kareyi dağılma özelliğiyle açınız.)
\item \(x=\sqrt{30+\sqrt{30+\sqrt{30+\cdots}}}\) olduğuna göre \(x\) kaçtır?
\end{enumerate}

\needspace{14\baselineskip}
\subsection*{Cevap Anahtarı}
{\small
\begin{multicols}{2}
\begin{enumerate}[label=\textbf{\arabic*.},leftmargin=*,itemsep=1pt]
\item \(0{,}2+1{,}2-0{,}5=0{,}9\).
\item \(3+(-4)=-1\).
\item \(3\sqrt2\); \ \(5\sqrt3\); \ \(\sqrt{64\cdot 2}=8\sqrt2\).
\item \((3+5-7)\sqrt2=\sqrt2\).
\item \(\sqrt{36}=6\).
\item \((\sqrt9)^3-(\sqrt[5]{32})^2=27-4=23\).
\item \(\sqrt7>2\Rightarrow|2-\sqrt7|=\sqrt7-2\).
\item \(49<60<64\Rightarrow 7\) ile \(8\) arası.
\item \(\dfrac{6\sqrt3}{3}=2\sqrt3\).
\item \(\sqrt[6]{3^2}\cdot\sqrt[6]{3^3}=\sqrt[6]{243}\).
\item \(\sqrt{20},\ \sqrt{18},\ \sqrt{19}\Rightarrow b<c<a\).
\item Payda \(6-5=1\Rightarrow\sqrt6+\sqrt5\).
\item \((\sqrt3-1)+(\sqrt3+1)=2\sqrt3\).
\item \(5+1=6\), \(5\cdot 1=5\Rightarrow\sqrt5-1\).
\item \(x^2=2+\sqrt2+\sqrt2+1=3+2\sqrt2\); \ \(3+2\sqrt2-2\sqrt2-2=1\).
\item \(x(x-1)=30=6\cdot 5\Rightarrow x=6\).
\end{enumerate}
\end{multicols}}


\vspace{6pt}
\noindent\textcolor{anamavi}{\rule{\textwidth}{0.5pt}}\\
{\small Gelecek hafta: Çarpanlara ayırma ve özdeşlikler.}

\end{document}
